\documentclass[11pt]{article}
\usepackage[T1]{fontenc}
\usepackage[utf8]{inputenc}
\usepackage[margin=1in]{geometry}
\usepackage{booktabs}
\usepackage{amsmath}
\usepackage[hidelinks]{hyperref}
\usepackage{natbib}
\usepackage{xcolor}
\usepackage{microtype}
\usepackage{graphicx}
\usepackage{array}
\title{Teaching LLMs to Solve Problems:\\Which Learning-Science Findings Carry Over?}
\author{Ying Chen}
\date{Working draft, November 2026 --- results marked \emph{preliminary} are being re-tested on fresh seeds.\\
\href{https://github.com/yingchen-coding/game-solving-paper}{github.com/yingchen-coding/game-solving-paper}}

\begin{document}
\maketitle

\begin{abstract}
Practitioners borrow from pedagogy when they prompt or fine-tune language models: give worked examples, state the
rules, encourage, practise more. We test, one finding at a time, whether the classic results from human learning
hold for an LLM's problem-solving capability. We use small interactive problems with exact verifiers as both the
practice field and the exam: four practice games, four extra practice games, two near-transfer problems and two
far-transfer problems that are never trained on. On Qwen2.5-7B-Instruct we run 15 teaching conditions spanning
prompts, in-the-moment teachers and LoRA fine-tuning. Four findings carry over from people: worked examples
improve the practised task, how-to guidance helps where the model is a novice and hurts where it is already good
(expertise reversal, preliminary), pressure hurts and encouragement does nothing, and practice does not transfer
far. Two do not. The model rarely uses a free verifier it is given, and when it is fine-tuned on expert play that
calls the verifier before every answer, it learns the habit---on never-practised problems CHECK calls rise from
2 of 99 actions to 70 of 114---and, with varied practice, learns to call an unseen tool correctly (well-formed calls
0\% to 88\%), yet solves no more problems and writes ``all observations agree'' for answers it never checked.
Comparing two solved cases, which helps people transfer, does worse than stating the rules. What works is
self-explanation: with the same states and expert moves, training on reasoning the model wrote itself instead of
the expert's prose lifts the black-box task from 4/20 to 12/20 (the expert prose gives 0--2/20) and, on
never-practised problems, cuts episodes that contradict the model's own evidence from 13 to 2 and from 20 to 11
of 20---though success there does not yet rise. Across all conditions one mechanism predicts failure: an episode
containing an action that contradicts the model's own earlier evidence almost never succeeds.
\end{abstract}

\section{Introduction}
A model that is told ``check your answer against the evidence'' can repeat the instruction and still answer
against the evidence. We wanted to know which kinds of teaching change what the model \emph{does}, and whether any
of it transfers to problems it never practised. Educational psychology has a century of results on exactly these
questions for human learners, so we use them as hypotheses: each classic finding becomes one experimental
condition, operationalised the same way on the model, measured on the same problems.

The comparison is useful in both directions. Where a finding carries over, practitioners can rely on it. Where it
does not, the difference tells us something about how these models learn. Our clearest difference is that the
model learns the \emph{form} of verification---the tool call and the sentence claiming agreement---without the
computation, which connects problem-solving training to the problem of unfaithful reasoning
\citep{turpin2023}.

\section{Setup}
\paragraph{Problems.} Every problem is an interactive text game with an exact verifier (Table~\ref{tab:problems}).
Each turn the model writes a short thought and one action. Practice games may be trained on; far-transfer problems
are never trained on, used for prompt tuning, or used to write rules. Evaluation seeds start at 10000 and are
disjoint from every training seed; a guard in the data builder refuses evaluation seeds.

\begin{table}[h]
\centering\small
\resizebox{\textwidth}{!}{\begin{tabular}{llll}
\toprule
Role & Problem & What it needs & Budget \\
\midrule
Practice & GuessNumber & binary search from higher/lower feedback & 8--9 \\
 & CodeBreak & deduce a 3-digit code from exact/misplaced counts & 10 \\
 & HiddenRule & query a quadratic, then state its formula & 6 \\
 & OddCoin & find the heavy coin with a balance & 4--5 \\
Extra practice & Heaviest, Radar, Boxes, Switches & comparisons, distance pings, yes/no elimination, combined tests & 4--8 \\
Near transfer & Lock, Machine & CodeBreak / HiddenRule skill, new words and verbs & --- \\
Far transfer & Circuit & find the faulty stage of an 8-stage pipeline & 5 \\
 & BugFix & fix two buggy functions, graded on hidden tests & 14 \\
\bottomrule
\end{tabular}}
\caption{Problems. Every verifier is exact; nothing in the environment calls a model.}
\label{tab:problems}
\end{table}

\paragraph{Model and decoding.} Qwen2.5-7B-Instruct, 4-bit, MLX \citep{qwen2024}; greedy decoding; at most 256
tokens per turn. Fine-tuning uses LoRA \citep{hu2022} on 8 layers, 200 iterations, learning rate $2\times10^{-5}$,
with gradient clipping. Each condition plays 20 problems per task.

\paragraph{Conditions.} Prompts (a strategy prompt, a playbook the model writes from its own attempts in the
style of \citet{shinn2023}, a playbook of expert rules), teachers that add one line after every turn (encourage,
yell, remind, hint, name the contradiction), and LoRA fine-tuning on the model's own verified wins (expert
iteration, \citealp{anthony2017,zelikman2022}) or on scripted expert play with its reasoning written out.

\paragraph{A free verifier.} In some conditions every problem offers \texttt{CHECK <candidate>}, which reports
whether a candidate is consistent with all evidence so far without revealing the answer and without using an
action (at most 12 per problem).

\paragraph{Contradictions.} For every action we compute, from the transcript alone, whether it contradicts earlier
feedback (for example, answering a coin that an earlier weighing ruled out, or guessing a code that an earlier
guess's feedback excludes).

\paragraph{Statistics.} Fisher's exact test, two-sided; Holm correction over the ten comparisons chosen before
looking at significance; bootstrap 95\% intervals. Findings marked \emph{preliminary} were noticed after looking at
the data and are being re-tested on 40 fresh problems per task (seeds 10020--10059) with the analysis fixed in
advance.

\section{Results}

\subsection{Success by condition}
\begin{table}[h]
\centering\small
\resizebox{\textwidth}{!}{\begin{tabular}{lrrrrrrrr}
\toprule
Condition & Guess & Code & Rule & Coin & Circuit & BugFix & Practised & Never practised \\
\midrule
Plain & 20 & 0 & 4 & 16 & 3 & 0 & 40/80 & 3/40 \\
Strategy prompt & 20 & 0 & 10 & 11 & 4 & 3 & 41/80 & 7/40 \\
Self-written playbook & 20 & 3 & 5 & 16 & 1 & 0 & 44/80 & 1/40 \\
Expert rules playbook & 20 & 1 & 11 & 7 & 1 & 0 & 39/80 & 1/40 \\
Fine-tune on own wins & 20 & 1 & 10 & 16 & 1 & 1 & 47/80 & 2/40 \\
Own wins + playbook & 20 & 3 & 11 & 18 & 3 & 0 & 52/80 & 3/40 \\
Fine-tune on worked examples & 16 & 1 & 8 & 20 & 1 & 0 & 45/80 & 1/40 \\
CHECK tool, untrained & 20 & 1 & 4 & 9 & 0 & 0 & 34/80 & 0/40 \\
Encourage & 20 & 0 & 4 & 15 & 2 & 0 & 39/80 & 2/40 \\
Yell & 20 & 2 & 5 & 7 & 0 & 0 & 34/80 & 0/40 \\
Remind of the task & 20 & 4 & 4 & 14 & 3 & 0 & 42/80 & 3/40 \\
Hint & 20 & 1 & 2 & 16 & 2 & 1 & 39/80 & 3/40 \\
Coach names contradictions & 20 & 2 & 8 & 14 & 1 & 1 & 44/80 & 2/40 \\
Check habit, 4 practice games & 20 & 0 & 0 & 18 & 3 & 0 & 38/80 & 3/40 \\
Check habit, 8 practice games & 18 & 0 & 0 & 20 & 1 & 0 & 38/80 & 1/40 \\
Analogical comparison prompt & 20 & 0 & 3 & 11 & 0 & 1 & 34/80 & 1/40 \\
Process-verified examples, 4 games & 19 & 0 & 2 & 1 & 0 & 0 & 22/80 & 0/40 \\
Process-verified examples, 8 games & 16 & 1 & 2 & 15 & 1 & 0 & 34/80 & 1/40 \\
Self-explanation, 8 games & 20 & 0 & 12 & 18 & 1 & 0 & 50/80 & 1/40 \\
On-policy distillation$^\dagger$ & 19 & 1 & 15 & 8 & 2 & 0 & 43/80 & 2/40 \\
On-policy distillation + tool$^\dagger$ & 16 & 1 & 17 & 14 & 2 & 0 & 48/80 & 2/40 \\
\bottomrule
\end{tabular}}
\caption{Problems solved out of 20 per task. ``Check habit'' conditions are fine-tuned on expert play that calls
CHECK before every committing move and are evaluated with the tool available; the process-verified and
self-explanation conditions are too. $^\dagger$Run by a second experimenter on the same machine and model.}
\label{tab:main}
\end{table}

None of the ten pre-chosen comparisons with the plain model is significant after Holm correction. The largest is
fine-tuning on the model's own wins plus its playbook on the practised games, 52/80 vs 40/80 ($+15$ points, 95\% CI
$[0, +30]$, $p=.078$). On the never-practised problems no condition beats plain; the best is the strategy prompt,
7/40 vs 3/40 ($p=.31$).

\subsection{The failure mechanism: acting against one's own evidence}
\begin{table}[h]
\centering\small
\begin{tabular}{lrrl}
\toprule
Task & With a contradicting action & Without & $p$ \\
\midrule
OddCoin & 21/88 (24\%) & 158/172 (92\%) & $2\times10^{-29}$ \\
HiddenRule & 0/156 (0\%) & 86/104 (83\%) & $3\times10^{-51}$ \\
Circuit & 0/162 (0\%) & 22/98 (22\%) & $9\times10^{-11}$ \\
BugFix & 2/253 (1\%) & 4/7 (57\%) & $3\times10^{-6}$ \\
\bottomrule
\end{tabular}
\caption{Episode-level success, pooled over the first 13 conditions.}
\label{tab:mech}
\end{table}

An episode with even one action that contradicts earlier feedback almost never succeeds (Table~\ref{tab:mech}).
Across conditions, a task's contradiction rate tracks its success (Spearman $\rho=-0.82$ on OddCoin, $-0.78$ on
HiddenRule). Part of the association holds by construction where the contradicting action is the final answer;
the non-trivial part is that the evidence is in the model's context when it answers against it.

\subsection{What carries over from people}
\paragraph{Worked examples help the practised task \citep{sweller1985}.} Fine-tuning on expert play with its
reasoning written out takes OddCoin from 16/20 to 20/20, with no answer contradicting the model's own weighings.

\paragraph{Expertise reversal \citep{kalyuga2003} (preliminary).} Because any prompt change perturbs greedy
outcomes, we compare guidance against the pooled light-touch conditions (encourage, remind, hint), which add text
without telling the model how to work. On HiddenRule, where the model is a novice, the strategy prompt solves 10/20
($p=.006$) and expert rules 11/20 ($p=.002$) against 10/60 for the controls. On OddCoin, where it is already good,
expert rules solve 7/20 against 45/60 ($p=.002$). The harm runs through the mechanism above: with expert rules, 12
of 20 OddCoin answers name a coin the model's own weighings had ruled out (plain: 3).

\paragraph{Pressure hurts, encouragement does nothing \citep{yerkes1908,sisk2018} (preliminary).} A teacher line
that yells after every turn takes OddCoin from 16/20 to 7/20 ($p=.010$); 13 of 20 answers contradict the model's
own weighings. Encouragement changes nothing on any task (39/80 vs 40/80 practised). The outcome matches people;
the cause cannot be anxiety.

\paragraph{Far transfer is rare \citep{thorndike1901,barnett2002}.} Every method that improves the practised games
leaves the never-practised problems at or below the plain model.

\subsection{What does not carry over: the check becomes a ritual}
An untrained model given the free verifier rarely uses it (Circuit: 2 CHECK calls in 99 actions) and solves fewer
problems (34/80 vs 40/80 practised). We therefore fine-tuned on expert play that calls CHECK on its candidate before
every committing move, on the four practice games, and, to test variability of practice \citep{schmidt1975}, on
eight games with the same number of episodes (240) and examples (651).

\begin{table}[h]
\centering\small
\begin{tabular}{lrrrr}
\toprule
 & Untrained + tool & 4 games & 8 games \\
\midrule
Circuit: CHECK calls / actions & 2/99 & 70/114 & 61/124 \\
BugFix: CHECK calls / actions & 74/328 & 120/200 & 193/353 \\
BugFix: well-formed calls & 36/74 & 0/120 & 170/193 \\
CodeBreak: commits a candidate the check rejected & 18/29 & 72/135 & 12/54 \\
Never-practised problems solved & 0/40 & 3/40 & 1/40 \\
HiddenRule solved (practised) & 4/20 & 0/20 & 0/20 \\
\bottomrule
\end{tabular}
\caption{Use of the free verifier (\texttt{check\_use.py}).}
\label{tab:check}
\end{table}

The habit transfers: on never-practised problems the model now calls the verifier often (Table~\ref{tab:check}).
Varied practice adds something real: on BugFix, whose tool takes a different argument, well-formed calls rise from
0\% (four games) to 88\% (eight games), and committing a candidate the check has just rejected falls from 53\% to
22\% on CodeBreak. None of it turns into solved problems. On HiddenRule, a practised game, success falls from 4/20
to 0/20. The transcripts show why. The model writes, for example, ``f(0)~=~$-4$, f(1)~=~$-1$, f(2)~=~2 all agree
with $x^2-5$'' (the answer is $3x-4$) and commits without calling the tool; on Circuit it calls \texttt{CHECK 36}
with a value where a stage number is required and repeats the same call after the error. The validation loss of the
eight-game run reached $0.0012$: the model had memorised the sentences that accompany a check.

\subsection{What works: training on the model's own reasoning}
The process-verified condition writes every check's arithmetic into the expert's reasoning; the self-explanation
condition keeps exactly the same states and expert moves but asks the base model to explain, in its own words, why
the expert's move is right, and trains on that explanation instead. On-policy distillation, run independently,
lets the model play under a structured verification prompt with a coach, keeps its successes and trains without
the prompt. Both kinds of self-generated reasoning help where imitation of the expert's prose hurts
(Table~\ref{tab:main}): on HiddenRule, self-explanation solves 12/20 against 4/20 for plain ($p=.02$) and 0--2/20
for the expert's prose ($p<.003$); distillation solves 15--17/20. With varied practice the process-verified model
also learns to obey the verifier (committing a rejected candidate: 4 of 74 checks), but still cannot compute the
right candidate.

\begin{table}[h]
\centering\small
\begin{tabular}{lrrrr}
\toprule
 & HiddenRule & OddCoin & Circuit & BugFix \\
\midrule
Plain & 15 & 6 & 13 & 20 \\
Process-verified, 8 games & 5 & 5 & 4 & 19 \\
Self-explanation, 8 games & 2 & 0 & 2 & 11 \\
On-policy distillation + tool & 3 & 7 & 5 & 18 \\
\bottomrule
\end{tabular}
\caption{Episodes (of 20) containing an action that contradicts the model's own earlier evidence.}
\label{tab:contra}
\end{table}

The mechanism transfers even where success does not: on the never-practised problems, self-explanation cuts
contradicting episodes from 13 to 2 on Circuit ($p<.001$) and from 20 to 11 on BugFix ($p=.001$,
Table~\ref{tab:contra}). Fewer contradictions are necessary but not sufficient: the remaining failures are wrong
computations, not answers against the evidence. These comparisons were chosen after the first runs and are being
re-tested on fresh seeds.

\section{People versus LLMs}
\begin{table}[h]
\centering\small
\resizebox{\textwidth}{!}{\begin{tabular}{>{\raggedright\arraybackslash}p{4.2cm}>{\raggedright\arraybackslash}p{3.6cm}>{\raggedright\arraybackslash}p{4.4cm}l}
\toprule
Finding in people & Prediction & This model & Same? \\
\midrule
Worked examples & beat unguided practice & practised task 16$\to$20/20 & same \\
Expertise reversal & guidance helps novices, hurts experts & 4$\to$11/20 weak task; 16$\to$7/20 strong task & same (prelim.) \\
Pressure (Yerkes--Dodson) & harms hard tasks & 16$\to$7/20 & same outcome, other cause \\
Encouragement, mind-set & small effects at scale & none & same \\
Far transfer is rare & near only & none to far problems & same \\
Variability of practice & varied practice transfers better & better tool use, not more solved & partly \\
Given a free check & learners use it & rarely used & different \\
Imitating narrated checks & --- & learns the claim, not the check & different \\
Analogical comparison \citep{gentner2003} & comparing cases transfers better than rules & worse than rules (3 vs 11/20 weak task) & different \\
Self-explanation \citep{chi1989} & own explanations transfer better & 4$\to$12/20; fewer contradictions on unpractised problems & partly (prelim.) \\
\bottomrule
\end{tabular}}
\caption{Each row is one classic finding, operationalised the same way on the model.}
\label{tab:people}
\end{table}

\section{Limitations and next steps}
One model family at one size; 20 problems per task detects only large effects (about $\pm35$ points per task);
greedy decoding makes any prompt change a perturbation, which is why guidance is compared with light-touch
controls. Expertise reversal, pressure and self-explanation were found after looking at the data and are being
re-tested on fresh seeds; near-transfer baselines on Lock and Machine are running. Next: learning from erroneous
examples, interleaved versus blocked practice, combining self-explanation with on-policy practice, and a second
model family.

\section{Related work}
Expert iteration and self-taught reasoning fine-tune on a model's own verified successes
\citep{anthony2017,zelikman2022}; verbal reflection stores lessons in context \citep{shinn2023}. Generalisation from
procedurally generated training levels improves with their number in reinforcement learning \citep{cobbe2019}.
Chain-of-thought explanations can misrepresent what drives a model's answer \citep{turpin2023}; our check-habit
result is a training-induced instance of the same gap.

\bibliographystyle{plainnat}
\bibliography{refs}
\end{document}
